\section{Method}
\label{sec:method}

\subsection{Predictors}
A shared ImageNet-pretrained EfficientNet-B0 encoder maps a face crop to a feature
vector, followed by one of three heads predicting VA in $[-1,1]^2$. The
\textbf{deterministic} head is trained with MSE and has no uncertainty. The
\textbf{Gaussian} head predicts a mean $\hat\mu_d$ and variance $\hat\sigma^2_d$
per axis $d\in\{v,a\}$ and is trained with the Gaussian NLL. The
\textbf{evidential} head~\cite{amini2020evidential} predicts
Normal--Inverse-Gamma parameters $(\gamma,\nu,\alpha,\beta)$ per axis, with
predictive mean $\gamma$ and predictive variance
$\beta(1+\nu)/(\nu(\alpha-1))$, the sum of its aleatoric term $\beta/(\alpha-1)$
and epistemic term $\beta/(\nu(\alpha-1))$. Its predictive distribution is a
Student-$t$ with $2\alpha$ degrees of freedom, which we use for its NLL.

From these we form two further methods without additional training.
\textbf{MC-dropout}~\cite{gal2016dropout} keeps the encoder's dropout ($p=0.2$,
the training rate) active for $20$ passes of the evidential model; its variance is
the mean predicted variance plus the variance of the pass means.
\textbf{Deep ensembles}~\cite{lakshminarayanan2017simple} combine the five
independently seeded models of one head, moment-matched in the same way. All
methods thus output a mean $\hat\mu(x)\in\mathbb{R}^2$ and a scale
$\hat\sigma(x)\in\mathbb{R}^2_{>0}$; for the deterministic head we set
$\hat\sigma\equiv 1$, which reduces every construction below to plain
absolute-residual conformal prediction.

\subsection{Per-axis conformal intervals (the common practice)}
Given a calibration set $\{(x_i,y_i)\}_{i=1}^n$ exchangeable with the test data,
the per-axis construction computes normalised scores
$s_{i,d}=|y_{i,d}-\hat\mu_d(x_i)|/\hat\sigma_d(x_i)$, takes
$q_d$ as the $\lceil(n{+}1)(1-\alpha)\rceil$-th smallest $s_{i,d}$, and outputs
$C_d(x)=\hat\mu_d(x)\pm q_d\,\hat\sigma_d(x)$. Each axis is covered with
probability at least $1-\alpha$. The VA point is covered when \emph{both} are,
which the union bound only guarantees with probability $1-2\alpha$.

\subsection{Joint conformal region}
We use one score per sample, the larger of the two normalised residuals,
\begin{equation}
  s_i=\max_{d\in\{v,a\}}\frac{|y_{i,d}-\hat\mu_d(x_i)|}{\hat\sigma_d(x_i)},
  \label{eq:maxscore}
\end{equation}
take $q$ as its $\lceil(n{+}1)(1-\alpha)\rceil$-th smallest value, and output the
rectangle $C(x)=\prod_{d}[\hat\mu_d(x)\pm q\,\hat\sigma_d(x)]$.

\begin{proposition}\label{prop:joint}
If the calibration and test pairs are exchangeable, then
$\Pr\{y\in C(x)\}\ge 1-\alpha$.
\end{proposition}
\begin{proof}
$y\in C(x)$ if and only if $s(x,y)\le q$, where $s$ is the scalar score of
Eq.~\eqref{eq:maxscore}. The standard split-conformal argument for a scalar score
then applies unchanged~\cite{lei2018distribution}.
\end{proof}

The \textbf{Bonferroni} alternative runs the per-axis construction at level
$\alpha/2$. It has the same guarantee but ignores that the two normalised
residuals are dependent, so it is conservative. The max-score adapts to that
dependence without modelling it. Copula
methods~\cite{messoudi2021copula} instead choose a per-axis level $t$ so that
the two calibration ranks fall below $t$ together with probability $1-\alpha$,
assuming independent axes ($t=\sqrt{1-\alpha}$), a Gaussian copula, or the
empirical copula of the calibration ranks; we compare against all three in
Section~\ref{sec:copula}. The empirical copula is the closest relative: with
ranks in place of $\hat\sigma$-normalised residuals it is a max-score region,
but since the ranks are computed on the calibration set itself its guarantee is
approximate rather than exact. We also compare with the ellipsoidal region of
Messoudi et al.~\cite{messoudi2022ellipsoidal}, a Mahalanobis score on the signed
normalised residuals; we estimate its covariance on one half of the calibration
split and its threshold on the other, so its guarantee is exact. An ellipse
does not factor into one interval per axis, so it cannot report a range for
valence and a range for arousal, which the rectangles can. The region inherits the predictor's
heteroscedasticity through $\hat\sigma$, so uncertain inputs get larger regions,
and it is the same object our on-device demonstrator uses as its abstention
signal.\footnote{Live in-browser demo, with the region drawn on the VA plane:
\demourl. No frames leave the device.}

\subsection{Group coverage and unseen expressions}
Marginal coverage can hide groups that are under-covered~\cite{romano2020malice}.
We report joint coverage separately for each AffectNet expression category and
summarise it by the worst seen category. The categorical labels are used only for
this audit; no model sees them. For inputs whose kind is absent from the
calibration data, no conformal method can promise coverage, so the remaining
safeguard is the uncertainty itself. We measure whether it rises on such inputs by
the AUROC of the total predictive standard deviation
$\sum_d\hat\sigma_d$ for separating them from the seen test inputs.

\subsection{Recalibrating for a new expression}
\label{sec:recal-method}
If $k$ labelled examples of the new expression can be collected, coverage on it
can be restored without retraining. \textbf{Group} recalibration (Mondrian
conformal~\cite{vovk2012conditional}) computes a separate max-score threshold
from those $k$ examples; by Proposition~\ref{prop:joint} applied within the
group, it has joint coverage at least $1-\alpha$ on that expression whenever
$k\ge\lceil 1/\alpha\rceil-1$ ($9$ at $\alpha=0.1$). Because an unseen
expression's error is largely a systematic offset (Section~\ref{sec:loeo}), we
also test \textbf{shift} recalibration: half of the $k$ examples estimate the
mean residual $\bar r$ of the expression, which is subtracted from its
predictions, and the other half calibrate the max-score threshold on the
corrected residuals $|y-\hat\mu-\bar r|/\hat\sigma$. The two halves are
disjoint, so this is again an exact split-conformal guarantee. Neither variant
changes the region of any other expression. Both require knowing which
expression a face belongs to at test time; we return to this limitation in
Section~\ref{sec:recal}.
